% Compilar desde la raíz con scripts/latexmk-build.sh.
% Producto: portada + UNA página de infografía. No es una entrega en Moodle.
% Datos pendientes deliberados: no inventar control ni semestre.
\documentclass[10pt,spanish]{article}
\usepackage[T1]{fontenc}
\usepackage[utf8]{inputenc}
\usepackage[spanish,es-nodecimaldot]{babel}
\usepackage[a4paper,landscape,margin=8mm]{geometry}
\usepackage{lmodern}
\usepackage{helvet}
\renewcommand{\familydefault}{\sfdefault}
\usepackage{graphicx}
\usepackage{tikz}
\usetikzlibrary{calc}
\usepackage{microtype}
\usepackage{hyperref}
\hypersetup{hidelinks,
  pdftitle={Tarea 4. Principales errores en la Matriz de Consistencia y el Círculo de Covey},
  pdfauthor={Martín Jonathan de la Cruz Muñoz},
  pdfsubject={Seminario I, Maestría en Gestión Administrativa}}
\pagestyle{empty}
\setlength{\parindent}{0pt}
\definecolor{navy}{HTML}{142D40}
\definecolor{wine}{HTML}{9F2539}
\definecolor{teal}{HTML}{087E83}
\definecolor{ink}{HTML}{203342}
\definecolor{muted}{HTML}{52636E}
\definecolor{paper}{HTML}{F2F5F7}
\definecolor{linegray}{HTML}{D9E2E8}
\definecolor{amber}{HTML}{8A5400}
\color{ink}

\newcommand{\studentname}{Martín Jonathan de la Cruz Muñoz}
\newcommand{\studentcontrol}{Por confirmar}
\newcommand{\studentsemester}{Por confirmar}
\newcommand{\scheduleddate}{6 de septiembre de 2026}

% Coordenadas en milímetros desde la esquina superior izquierda.
\newenvironment{canvas}{%
  \begin{tikzpicture}[remember picture,overlay]
  \begin{scope}[shift={(current page.north west)},x=1mm,y=-1mm]
}{\end{scope}\end{tikzpicture}}

\newcommand{\errorcard}[8]{%
  \begin{scope}[shift={(#2,#3)}]
    \filldraw[fill=white,draw=linegray,rounded corners=2mm]
      (0,0) rectangle (90,43);
    \fill[#4,rounded corners=1mm] (3,3) rectangle (11,11);
    \node[text=white,font=\bfseries\fontsize{11}{12}\selectfont]
      at (7,7) {#1};
    \node[anchor=north west,text width=73mm,inner sep=0pt,
      text=#4,font=\bfseries\fontsize{10.3}{11.3}\selectfont]
      at (14,3) {#5};
    \node[anchor=north west,text width=83mm,inner sep=0pt,
      align=left,font=\fontsize{9.1}{10.6}\selectfont]
      at (3.5,13) {%
      \textbf{Falla:} #6\par\vspace{0.8mm}
      \textbf{Ejemplo:} #7\par\vspace{0.8mm}
      \textcolor{teal}{\textbf{Evita:}} #8};
  \end{scope}}

\begin{document}
\null
\begin{canvas}
  \fill[paper] (0,0) rectangle (297,210);
  \fill[navy] (0,0) rectangle (88,210);
  \fill[wine] (88,0) rectangle (91,210);
  \node[anchor=north west,inner sep=0pt] at (12,13)
    {\includegraphics[width=62mm]{ITESCA/assets-itesca/web/logo-itesca-contorno.png}};
  \node[anchor=north west,text width=63mm,text=white,inner sep=0pt,
    font=\bfseries\fontsize{12}{15}\selectfont] at (12,57)
    {INSTITUTO TECNOLÓGICO\\SUPERIOR DE CAJEME};
  \node[anchor=north west,text width=63mm,text=white,inner sep=0pt,
    font=\fontsize{11}{15}\selectfont] at (12,83)
    {Posgrado e Investigación\\[4mm]Maestría en Gestión Administrativa\\[4mm]Seminario I};
  \node[anchor=north west,text width=63mm,text=white,inner sep=0pt,
    font=\fontsize{10}{13}\selectfont] at (12,167)
    {Trabajo individual\\Infografía académica\\[3mm]Cajeme, Sonora};
  \node[anchor=north west,text=wine,inner sep=0pt,
    font=\bfseries\fontsize{12}{14}\selectfont] at (104,16)
    {TAREA 4 \quad / \quad UNIDAD 2};
  \node[anchor=north west,text width=179mm,text=navy,inner sep=0pt,
    font=\bfseries\fontsize{25}{29}\selectfont] at (104,28)
    {Principales errores en la\\Matriz de Consistencia\\y el Círculo de Covey};
  \node[anchor=north west,text width=176mm,text=muted,inner sep=0pt,
    font=\fontsize{11.5}{15}\selectfont] at (104,66)
    {Reconocer errores, alinear el proyecto y actuar sobre lo que puede mejorarse.};
  \draw[linegray,line width=0.7pt] (104,85) -- (283,85);
  \node[anchor=north west,inner sep=0pt,font=\fontsize{11}{15}\selectfont]
    at (104,93) {\renewcommand{\arraystretch}{1.18}%
    \begin{tabular}{@{}p{40mm}p{132mm}@{}}
      \textbf{Nombre completo} & \studentname \\
      \textbf{Carrera} & Maestría en Gestión Administrativa \\
      \textbf{Semestre} & \studentsemester \\
      \textbf{Número de control} & \studentcontrol \\
      \textbf{Docente} & Dra. Carla Olimpya Zapuche Moreno \\
      \textbf{Fecha de entrega} & \scheduleddate\ (según calendario local) \\
      \textbf{Periodo} & Agosto--noviembre de 2026
    \end{tabular}};
  \node[anchor=north west,text width=178mm,text=amber,inner sep=0pt,
    font=\fontsize{9.5}{12}\selectfont] at (104,166)
    {\textbf{Versión para revisión.} Completar número de control y semestre;
    confirmar la fecha vigente antes de entregar.};
  \node[anchor=north west,text width=177mm,text=muted,inner sep=0pt,
    font=\fontsize{9}{12}\selectfont] at (104,187)
    {Elaboración: 12 de septiembre de 2026.\\
    Estructura: portada + una página de infografía.};
\end{canvas}
\newpage
\null
\begin{canvas}
  \fill[paper] (0,0) rectangle (297,210);
  \fill[navy] (0,0) rectangle (297,31);
  \fill[wine] (0,31) rectangle (297,32.2);
  \node[anchor=north west,text=white,inner sep=0pt,
    font=\bfseries\fontsize{20}{23}\selectfont] at (8,5)
    {8 errores que rompen la coherencia};
  \node[anchor=north west,text=white,inner sep=0pt,
    font=\fontsize{10.5}{13}\selectfont] at (8,17)
    {Matriz de Consistencia + Círculo de Covey \quad | \quad Detectar $\rightarrow$ corregir $\rightarrow$ verificar};
  \node[anchor=north east,text=white,inner sep=0pt,align=right,
    font=\bfseries\fontsize{9}{12}\selectfont] at (289,6)
    {SEMINARIO I\\MGA / TAREA 4};
  \node[anchor=north west,text width=281mm,inner sep=0pt,
    font=\fontsize{9.2}{11}\selectfont] at (8,35)
    {\textbf{Ruta de consistencia:} problema $\rightarrow$ pregunta $\rightarrow$ objetivos
    $\rightarrow$ hipótesis, si procede $\rightarrow$ variables/categorías $\rightarrow$ método $\rightarrow$ evidencia.};

  \errorcard{1}{8}{46}{wine}{Confundir tema con problema}
    {Nombrar un campo no identifica un vacío de conocimiento.}
    {«Digitalización» es tema; falta saber cómo se relaciona con la recompra.}
    {Expón la brecha y susténtala con antecedentes.}
  \errorcard{2}{103}{46}{wine}{Problemas demasiado generales}
    {Sin límites, no se sabe a quién ni qué estudiar.}
    {«¿Cómo mejorar todas las empresas?» es inabarcable.}
    {Acota fenómeno, población, lugar y periodo: uso digital y recompra; comercios de Cajeme, 2026.}
  \errorcard{3}{198}{46}{wine}{Objetivos vagos o no medibles}
    {«Conocer mejor» no precisa el resultado observable.}
    {«Conocer la digitalización» no indica qué analizar.}
    {Usa verbo + variables + límites: analizar la asociación entre uso digital y tasa de recompra en esos comercios.}

  \errorcard{4}{8}{93}{wine}{Hipótesis desconectadas}
    {La respuesta tentativa cambia las variables de la pregunta.}
    {Se pregunta por uso digital y recompra, pero se plantea una hipótesis sobre motivación laboral.}
    {Propón una relación contrastable entre las mismas variables, población y periodo.}
  \errorcard{5}{103}{93}{wine}{Método incongruente}
    {Una encuesta transversal no basta para atribuir causalidad.}
    {«Demostrar que digitalizar causa más recompra».}
    {Ajusta el objetivo a asociación o justifica un diseño causal; define muestra, medición y análisis acordes.}
  \errorcard{6}{198}{93}{teal}{Variables mal ubicadas en Covey}
    {Clasificar por importancia confunde preocupación con influencia.}
    {Ubicar la inflación como algo que el tesista puede controlar.}
    {Ubica factores según tu capacidad real de actuar. No confundas influencia con variable independiente.}

  \errorcard{7}{8}{140}{wine}{Repetir ideas entre apartados}
    {Copiar la misma frase oculta la función de cada apartado.}
    {«Mejorar ventas» en problema, objetivo y método.}
    {Diferencia: brecha, acción y procedimiento. Mantén los mismos conceptos, no párrafos duplicados.}
  \errorcard{8}{103}{140}{teal}{Confundir influencia con control}
    {Poder solicitar apoyo no garantiza la conducta ajena.}
    {«El gerente entregará todos los datos» no depende solo del tesista.}
    {Controla tu solicitud; negocia el acceso, respeta la confidencialidad y prevé una alternativa.}

  % Los dos círculos de Covey son anidados, no conjuntos disjuntos.
  \filldraw[fill=white,draw=linegray,rounded corners=2mm]
    (198,140) rectangle (288,183);
  \node[anchor=north west,text=teal,inner sep=0pt,
    font=\bfseries\fontsize{10}{12}\selectfont] at (201.5,143)
    {Covey: enfoca tu siguiente acción};
  \filldraw[fill=paper,draw=muted] (219,166) circle[radius=14mm];
  \node[font=\bfseries\fontsize{7.3}{8}\selectfont] at (219,155.5)
    {Preocupación};
  \filldraw[fill=teal!12,draw=teal,line width=0.7pt]
    (219,170) circle[radius=8mm];
  \node[align=center,text=teal,font=\bfseries\fontsize{7.2}{8}\selectfont]
    at (219,169) {Influencia\\Actuar};
  \node[anchor=north west,text width=48mm,inner sep=0pt,
    font=\fontsize{8.6}{10.1}\selectfont] at (237,152)
    {\textbf{Directo:} revisar objetivos.\\[1mm]
     \textbf{Indirecto:} negociar acceso.\\[1mm]
     \textbf{Sin control:} inflación; prever escenarios.};

  \node[anchor=north west,text width=281mm,inner sep=0pt,
    font=\fontsize{8.7}{10.4}\selectfont] at (8,186)
    {\textbf{Antes de entregar:} cada objetivo debe tener evidencia y método.
    Las hipótesis dependen del enfoque y alcance; Covey prioriza acciones, no sustituye el diseño científico.
    \textbf{Ejemplos hipotéticos:} no son resultados de una investigación.};
  \draw[linegray] (8,195) -- (289,195);
  \node[anchor=north west,text width=281mm,text=muted,inner sep=0pt,
    font=\fontsize{7.4}{8.8}\selectfont] at (8,197)
    {\textbf{Fuentes:} Hernández Sampieri, R., Fernández Collado, C. y Baptista Lucio, P. (2014).
    \textit{Metodología de la investigación} (6.ª ed.). McGraw-Hill; pp. 37--39, 104--109 y 152--153.\\
    Covey, S. R. (2003). \textit{Los 7 hábitos de la gente altamente efectiva}. Paidós; primer hábito, pp. 50--52 del PDF consultado.
    \hfill Elaboración: \studentname.};
\end{canvas}
\end{document}